\documentclass[11pt]{article}
\usepackage{amsmath,amssymb,amsthm}
\usepackage[margin=1in]{geometry}

\newtheorem{theorem}{Theorem}
\newtheorem{lemma}[theorem]{Lemma}

\title{Problem 561: Divisor Pairs}
\author{Project Euler Solutions}
\date{}

\begin{document}
\maketitle

\section{Problem Statement}

Let $\mathrm{Pairs}(n)$ denote the number of ordered pairs $(d_1,d_2)$ of positive divisors of~$n$ with $d_1 \le d_2$ and $n = d_1 \cdot d_2$. Define
\[
S(N) = \sum_{n=1}^{N}\mathrm{Pairs}(n).
\]
We seek $S(10^{10})$.

\section{Mathematical Foundation}

\begin{theorem}[Reduction to Divisor Counting]
For every positive integer~$n$,
\[
\mathrm{Pairs}(n) = \left\lceil \frac{\tau(n)}{2}\right\rceil = \frac{\tau(n) + [n \text{ is a perfect square}]}{2}.
\]
\end{theorem}

\begin{proof}
The divisors of~$n$ partition into pairs $(d, n/d)$ with $d < n/d$, plus a singleton $\{d\}$ when $d^2 = n$. Each pair yields one ordered pair $(d_1,d_2)$ with $d_1 < d_2$; the singleton yields $(d,d)$. Hence
$\mathrm{Pairs}(n) = \lfloor\tau(n)/2\rfloor + [n \text{ is a perfect square}] = \lceil\tau(n)/2\rceil$.
\end{proof}

\begin{lemma}[Summatory Formula]
\[
S(N) = \frac{1}{2}\!\left(\sum_{n=1}^{N}\tau(n) + \lfloor\sqrt{N}\rfloor\right).
\]
\end{lemma}

\begin{proof}
Summing over $n$, the Iverson bracket $[n\text{ is a perfect square}]$ contributes $\lfloor\sqrt{N}\rfloor$.
\end{proof}

\begin{theorem}[Hyperbola Method]
\[
\sum_{n=1}^{N}\tau(n) = 2\sum_{d=1}^{\lfloor\sqrt{N}\rfloor}\left\lfloor\frac{N}{d}\right\rfloor - \lfloor\sqrt{N}\rfloor^2.
\]
\end{theorem}

\begin{proof}
Write $\sum_{n=1}^{N}\tau(n) = \#\{(a,b): ab \le N\}$. The Dirichlet hyperbola method splits the lattice points below the hyperbola $ab = N$ into those with $a \le \sqrt{N}$, those with $b \le \sqrt{N}$, and subtracts the doubly-counted square $a,b \le \sqrt{N}$.
\end{proof}

\begin{theorem}[Sublinear Evaluation]
The sum $\sum_{d=1}^{M}\lfloor N/d\rfloor$ can be computed in $O(\sqrt{N})$ time by grouping consecutive~$d$ with the same quotient $\lfloor N/d\rfloor$.
\end{theorem}

\begin{proof}
There are $O(\sqrt{N})$ distinct values of $\lfloor N/d\rfloor$ for $d \in \{1,\ldots,N\}$: if $q \le \sqrt{N}$ or $d \le \sqrt{N}$. Each block of consecutive~$d$ sharing the same quotient is processed in $O(1)$.
\end{proof}

\section{Algorithm}

\begin{enumerate}
\item Compute $M = \lfloor\sqrt{N}\rfloor$.
\item Using the block decomposition, compute $T = \sum_{d=1}^{M}\lfloor N/d\rfloor$ in $O(\sqrt{N})$.
\item Set $\mathrm{sum\_tau} = 2T - M^2$.
\item Return $(\mathrm{sum\_tau} + M)/2$.
\end{enumerate}

\section{Complexity Analysis}

\begin{itemize}
\item \textbf{Time:} $O(\sqrt{N})$. The block decomposition visits $O(\sqrt{N})$ quotient classes.
\item \textbf{Space:} $O(1)$.
\end{itemize}

\section{Answer}

\[
\boxed{452480999988235494}
\]

\end{document}
